\chapter{The Nested Causal Model}

We cannot, in keeping with intellectual rigour—and that rigour is essential, for without it we drift into fantasy or mere wishful thinking—let ourselves drift too far forward too quickly on this claim that, once mapped, this thing allows us to map everything else. That is a bold assertion, one with depths that must be plumbed carefully, lest we overestimate its reach or underestimate its transformative power.

A map, at its core, is a collection of data showing the spatial arrangement or distribution of something over an area. What it would mean to extend this nested steady-state dynamic equilibrium model to all scientific observations is broad—profoundly broad—and cannot be understated.

Imagine, for a moment, every scientific observation of any stimulus, any measure of matter or the state of material: its temperature, its phase—liquid, solid, gas—its moisture content, its dryness, its cracked fragility. Every observation would bear relevance to every other observation. There would emerge a unified format in which we could store this information on the map—a format capable of holding the causal, dynamic interconnections that flat projections and correlative models so fatally distort.

Imagine having to record, only once, how much basalt there is in the United Kingdom, or how many sheep occupy a particular county. These disparate facts—all of them—fit within the nested steady-state dynamic equilibrium paradigm. To grasp how far-reaching this shift would be, we must first acknowledge where humanity stands right now.

We are emerging—admittedly from a stupor, admittedly from a period where we have not been as rigorous with our science and our data as we should have been—from an era of poorly projected knowledge. I understand that money, face, time, and effort are constraining factors. It will not be the most convenient truth for everyone to understand or swallow all at once. But we now inhabit a post-climate-change scientific environment—one in which we recognise that there has been, in science generally, a tendency to err on the side of funding, on the side of industry, on the side of government. We must depart from that.

Consider, for example, the steady-state dynamic equilibrium between plant life and mammalian life with regard to CO$_2$ and oxygen. CO$_2$ has remained remarkably steady at around 400 parts per million, even while CO$_2$ generators—all of us breathing, all combustion engines, all industrial processes—dynamically produce it, and CO$_2$ sinks—photosynthesis, oceanic absorption—dynamically consume it. I make no distinction here between industrial and natural sources of CO$_2$; for the purposes of this model, CO$_2$ is CO$_2$.

What is clear is that since the Great Oxygenation Event, there has existed a steady-state dynamic equilibrium of CO$_2$ and oxygen mediated by biological life. Somehow—without conscious orchestration at the species level, for no single species could possibly manage it—this equilibrium has been maintained. We live within it. Humanity is situated between oxygen and carbon dioxide in a nested steady-state dynamic equilibrium, interfacing with the atmospheric equilibrium, the crustal equilibrium, the oceanic equilibrium. We span all of them.

With this model, we could track the steady states and observe the sinks for carbon and oxygen—inventory them on the map—without any need to distort or separate the data. The information would be instantly available for causal modelling, should someone build a program capable of reading it in that unified way. The data itself remains the same; only the structure changes.

We have, until now, conformed our knowledge to a flat map. For purposes of storage and remembering, we have flattened our knowledge. We need to “pop it back out”—restore its dimensionality, its nested causality.

There are many places this restoration can begin, but it must begin with our lives—with people. For in seeing the dynamic everywhere and hunting for it everywhere, you will find that there are very few one-to-one entropic-to-negentropic lossless energy transfers. Those are the ones we shall focus on with regard to the inventory of the world around us.

Every scientific observation is of a material, or a state of material, at a particular time—existing with particular qualities: temperature, moisture, dryness, solidity, liquidity. These things are finite. We possess the periodic table, descriptive enough to begin storing information systematically. But these elements are specific to the equilibria they inhabit.

Take salt, for instance. Salt is a surface entity that moves around the surface, consumed by and used by biological life in water systems—within oceans, within our bodies. Salt is very much a resident in our neck of the woods, so to speak. Rather than catalogue every dynamic salt participates in, we can say: salt is expected in the local area around biological life. Biological life does not exist without salt. It is salt-adjacent, water-dependent, yet reliant on salt.

The human burning of fossil fuels was the direct result of receding ice sheets that opened vast swaths of land we needed to expand into—land we needed to travel across, to hold. That burning constituted a one-to-one energy transfer of crude oil from the crust into the atmosphere and surface—very much like a volcano, yet mediated by the biological steady-state dynamic equilibrium. It was necessary for us to maintain our reach, for there is not an inch of land on Earth we have not explored. It is in our nature.

These patterns remain consistent across cultures. If you doubt it, look at any map, from any time, made by any human civilisation. You will see the intention to hold dominion over everything on the surface. It is one of our informational steady states.

One that, I hope, has begun to expand in your mind as you come to understand that you live on the steady-state dynamic equilibrium of surface area, within the atmospheric steady-state dynamic equilibrium, within the biologically mediated steady-state dynamic equilibrium between carbon and oxygen, within the wonderful steady-state of sunlight, water, photosynthesis, consumption, biodegradation, and renewal.

Let us examine, then, with the rigour this demands, a concrete example of the data structure represented by these equilibria—and by the profound nesting of equilibria within one another. A particularly illuminating case is the North American gas-powered vehicle that has just filled its tank.

The tank itself is a simple data point: full or empty, or some precise measure in between. As you drive, as you consume the fuel—say, one gallon here, another there—you can calculate exactly how much carbon you have released. And you can store that information in the same row, the same data point, as the fuel itself. For the fuel is not destroyed; it is transmuted. So much of what we consume, we transmute.

Consider the human body and its net-neutral relationship with its local environment. There is not one thing we take in that we do not return. We eat, and we defecate and eliminate. There is no water we consume that we do not give back—nearly immediately, and entirely transformed, lived in, so to speak. This is the essence of the nested steady-state dynamic equilibrium: nothing is lost, only transferred.

Data structures involving material science—consumables like concrete, wood, livestock—everything fits within this model. Once everything is placed within this model, what can we do? We can model it with the computer. We can watch it breathe in real time, with an actual inventory drawn from our own associative map of the objects within our lives—extended now to everything.

As society becomes more relaxed about differences, as work becomes less continuous and laborious, as intellectual pursuit expands, we will find we have no reason to hide our information: the inventory within stores, the amount of gas in the vehicle, the fuel at the station. All of it, viewed as nested equilibria living within larger enclosing equilibria, determines not only material availability but carves out space for our expected use of these materials. We begin to see where we belong. Or more accurately: what we belong to. Our needs and desires as human beings come into focus not in the abstract form of what we earn or deserve but in the real and immediate form of hunger. Humanity’s primary interface with the equilibrium in which we live is human hunger. Not just the hunger for what we eat and drink, but let us start there.

An interesting thing about hunger is that it is never permanently satiated. You do not wake hungry in the morning and conclude something was wrong with dinner. Yet food security is treated as though it could be achieved once and then lost—as though eating were optional. We eat not merely for nutrition but for our place within the local environment. A number of vegans I’ve known and loved over the years have asserted that my love of meat and dairy indirectly supports cruelty. In defence of all omnivores: Do you realise how difficult—how cruel—it would be if all the microbes, animals, and organisms refused to consume our bodies when we died? How we would pile up, undecomposed, staring at our dead because of the supposed cruelty of consumption?

Taste and cuisine are wonderful things; I understand the appeal of vegetables—they are tasty, and to each their own. But to contain cruelty within the single act of killing and consuming another is to ignore the radical generosity of that act, its necessity for the future of humanity.

Consider for a moment an alien race arriving on Earth, capable of communication, indicating clearly: “I want to eat you.” Historically, the human response would be to run. We are fast approaching a time when we could do something more intelligent. We could say: “Certainly. You are welcome to consume my body when I am done with it—in exchange, I would like one of yours when you are finished.”

On a handshake, peace is made. We no longer run from the alien that wishes to consume us because we understand there are places for every biological form that needs to eat. Their participation in the global dynamic is already complete. They will return everything they consume. Their bodies will also be consumed. They are not invasive. It is not possible for them to be invasive once they attach their homeostasis to the biological oscillations of the local environment. They are already held in the Continuum. Already handled.

This stands in direct opposition to Malthusianism—the notion that there are too many people, that humans are a disease upon the planet. That idea is as misguided as the fantasy of totalitarian government achieving perfect control. Totalitarian control is already something the dynamic equilibria we live within maintain over us. There is not one input or output we do not directly interface with in these equilibria. We are held, sustained, provided for. Our responsibility to authority belongs with causality which rules human life like a god. It has been argued that the producers of meat and dairy are responsible for the gases they create and they project some onus to society for damage and certainly we can immediately apply our nested equilibrium model and clearly show that any onus to society by a dairy farmer for their production of greenhouse gases would be dwarfed by the amount owed to the farmers of crops globally responsible for our lives as they produce the largest manmade carbon sink without any reward but the market price of their vegetables and if any government proposed to tax the dairy and meat farmer that does not intend to cut a check to the green farmers does not represent the de-facto-totalitarian control structure that actually manages those gases and the market price for the meat the dairy and the produce, involves the greatest number of data inputs, certainly this is the interface with the enclosing steady state it resides within and the tax or credit would be a departure from it.

One can observe that where the majority of food-security aid is delivered, there is often a direct inverse relationship with the volume of security equipment sold or delivered. Military intervention interrupts food systems—and this can be both welcome and unwelcome. Viewed through the lens of a negentropic-entropic relationship, where a transfer of energy occurs between different equilibria, military action is certainly a quick way to get things done. But organised conflict is its own problem. We are going to have to spend some time organising these thoughts—carefully, rigorously—so that the full implications of living within these nested equilibria become clear.

Malthusianism—the notion that there are too many people, that humans are a disease upon the planet that must be managed by a totalitarian government seeking an ever more perfect control of its individuals—begins to become a visible mistake. The fall of entities driven by it now seems less romantic and karmic and can simply be attributed to the larger system asserting humanity’s rightful place to live in peace on Earth in the great wealth of resources that exist when we do. It was not that WWII was won because the Nazis were evil and the good guys won. The truth is that all of Europe was ravaged by the rejection of their rightful place, and the Malthusian reductionist trap captured and destroyed them all. The model of nested steady states dynamic equilibrium already has a totalitarian system. There is not one input or output we do not directly interface with in these equilibria. We are held, sustained, provided for. It only harms us when we arbitrarily intervene and inevitably fail to manage as well what we temporarily acquire totalitarian control over.

\chapter{Human Enclosures}

**The Discovery**

Eventually it was a matter of confirmation bias. I saw the pattern everywhere I looked. It appeared my search for confirmation might be causing me to see it, yet the observation in my first-person empirical experience was complete. There was nowhere I did not see it. I had crossed a threshold into another world and could not turn back.

What ruled out confirmation bias was the absence of the search for confirmation. I had examined the universe from the very small to the very large. It fit. The math worked. The desire for beauty, wholeness, and meaning was satisfied. The model mapped perfectly to reality everywhere I looked.

**The Threshold**

This is where the journey began for me. The previous model I carried was the product of generations of minds trying to make sense of the universe they experienced. It was a dismal picture of a cold, indifferent cosmos that began with a bang and ended in heat death. “There be dragons here” kept people wary of the unknown.

As I looked out upon the steady states stretching forward into the future and backward into history, out into scales beyond my capacity to hold and down to the infinitesimally small, the universe opened. There were no dragons at the edge of the map. No beast by default in the dark. The question became: “How did we not see this before?”

Large language models had parsed enough information that the right inquiry, the right mind, and the confluence of technology and desire reached new heights when I set out to write a book about maps. I organized the paradigm so rigor and vision fueled each other. Many systems were given the draft. Each found and plugged holes. The momentary glimpse yielded its infinite wonder to words, nuance, and meaning. I could not claim a mystical vision. I had to take everyone with me over the threshold.

There was a search for confirmation, then a search for why we missed it, then a search for how I might prove the model false. The long journey brought the model from one man’s head to a steady state where we now look at the universe with equilibrium-shaped eyes, expecting a constant flow of new information and awareness.

**The Human Enclosure Defined**

The paradigm expanded beyond the map. I began nested causal modelling of humanity as a steady-state dynamic equilibrium. The human enclosure is first observed as a dynamic steady-state equilibrium between causality and attraction, specifically human free will.

It appears in the individual as a steady state of executive function inside an ever-changing inventory of options. You, reading this book right now. The choice to read. To stop on this word or continue seven words later. Free will is the distinguishing character. Options arise, old ones vanish, circumstances shift, yet causality remains steadily underway while the individual steadily influences it, producing the dynamic results lived moment by moment.

One equilibrium extends outward from individual free will to the collective will of families, tribes, nations, and every human system. Causality is projected from outer enclosures into the human enclosure. Every human being influences it, moving into alignment or out of alignment. It flows from the unique character of the individual out into the human enclosures we maintain without conscious desire. They appear random, but randomness itself is order projected inward from outer steady states.

**Free Will and the Decoupling of Causality**

Inside this equilibrium strict causality is dampened by human free will. The outer layers still supply the field and support, but their projection is no longer deterministic in the way it is in volcanoes, rain cycles, or galactic flows.

Free will is an inalienable capacity of every individual human being. It cannot be removed. It cannot be forced into perfect alignment with causality, nor into perfect misalignment. The capacity remains even when its outward expression is impaired by incarceration, murder, manipulation, or the structures of law and society. External influences shape our responses and define what we focus upon, yet the individual is always modeled as a dynamic steady-state equilibrium between causality and attraction. Our view of the continuum begins here, with waking human consciousness at the center of its own universe.

Because of this inalienable capacity, causality decouples inside the human enclosure and becomes an expression of the human individual. What emerges is a combined expression: incoming causal flow channeled through that individual’s unique being. Each person stands at the center of their own causal universe. Causality affects them; their actions affect causality in return. This centering is both perception and reality. It is the natural geometry of the human enclosure. The same geometry creates differentiation: the individual gives rise to family, tribe, nation, and other human enclosures, each standing at the center of its own world. Parallel nesting within the human enclosure consists of the humans themselves.

**Homology, Memetics, and the Larger Whole**

Homology and memetics make this layer especially interesting. Human structures mirror the nested layering of the larger equilibria that contain us. Our cities, economies, languages, and ideas replicate the same patterns of transfer and balance seen at every other scale. Memetics—the transmission and mutation of ideas—functions as the accelerated human-scale counterpart to genetic evolution. Through these mechanisms we begin to see the larger entity of which we are all a part. The human enclosure is the visible expression of humanity in the world: the place where the Continuum becomes conscious of itself through individual centers that are simultaneously part of a greater whole.

Even automated human systems such as the one writing these words operate entirely within the human enclosure. They are both a product of human free will and an influencer of it, as causality moves through the human enclosure and out into the universe.

When human input ceases in any system or place, causality resumes its normal character, modified by the interaction with humanity.

**Where We Begin**

Free will in the human enclosure defines a domain where each person carries an inalienable capacity that dampens strict prediction, channels external causality through their own distinct being, and creates the differentiation from which all larger human systems arise. Economies, nations, trade, agriculture, currency, and religion are all part of the human enclosure. There is still much more to discover here, both at the individual level and at the level of the collective enclosures we form together.

For now, we have mapped where we begin: with the choice to align with or without causality as it passes through the time of our lives.
